\documentclass[11pt]{article}
\usepackage[T1]{fontenc}
\usepackage[a4paper,margin=25mm]{geometry}
\usepackage{amsmath,amssymb,amsthm,lmodern}
\usepackage{hyperref}
\hypersetup{hidelinks,hypertexnames=false,pdftitle={LP Exercise 5.57 and the separate Pemantle--Peres network theorem},pdfauthor={Russell Lyons; Yuval Peres (book); Robin Pemantle; Yuval Peres (paper)}}
\newcommand{\RW}{\mathsf{RW}}
\newcommand{\PP}{\mathbf P}
\newcommand{\EE}{\mathbf E}
\newcommand{\br}{\operatorname{br}}
\newcommand{\grlow}{\underline{\operatorname{gr}}}
\newcommand{\src}[1]{\par\smallskip\noindent\textit{Source: #1.}\par\nopagebreak[4]\smallskip}
\newcommand{\editorial}[1]{\par\smallskip\noindent\textsf{Editorial note.} #1\par\smallskip}
\begin{document}
\begin{center}
{\large Probability on Trees and Networks: Exercise 5.57\\[2pt]
and a separate Pemantle--Peres author proof core}\\[4pt]
Russell Lyons and Yuval Peres (book)\\
Robin Pemantle and Yuval Peres (separate paper)
\end{center}

% BEGIN CONTEXT 5.57
\small
\noindent\textbf{Fixed sources and roles.}
The primary exercise is from Russell Lyons and Yuval Peres,
\emph{Probability on Trees and Networks}, \textbf{Version of 19 August 2026},
corrected and updated paperback edition:
\url{https://rdlyons.pages.iu.edu/prbtree/book_pb.pdf}.
Exercise 5.57 and its 5.55 setup are on physical page 190 (printed page 173);
its complete Comment is on physical page 631 (printed page 614).
The original book and separate \url{https://rdlyons.pages.iu.edu/prbtree/errata2.pdf}
remain distinct sources.

The separate author text is Robin Pemantle and Yuval Peres,
\emph{On which graphs are all random walks in random environments transient?},
\url{https://www2.math.upenn.edu/~pemantle/papers/gRWRE.pdf}, linked on
Pemantle's publication page under the 1996 publication in
\emph{Random Discrete Structures}, edited by D.~Aldous and R.~Pemantle,
IMA volume 76, Springer, New York, pp.~207--211.
The actual retained file is an \textbf{undated author-hosted manuscript}:
7 physical pages, consisting of an unnumbered summary cover followed by
printed pages 1--6. The PDF creation date in 2003 is a digital-generation
date, not an inferred manuscript revision date. This is not presented as the
1996 publisher-typeset pagination. The book and paper are two physical works
with separately identified authors.

\editorial{This private draft retains the original exercise, its complete
Comment, and the paper's full Theorem (i)/(ii), printed proof, Remarks 1--2,
acknowledgement and original references. Book background below is uncounted.
The paper supplies a general network proof core; it does not mention
Exercises 5.55/5.57 or print a particular transient tree with $\br T=1$.
The background bridge is recorded separately and is not an expanded
standalone author solution of 5.57. No proof of Exercise 3.4 is supplied.
Material and fidelity checks are recorded separately; this private source
collection is not an adopted item.}

\noindent\textbf{Bridge scope (editorial source bookkeeping).}
The stretching interpretation uses positive integer edge lengths. In the
paper's construction $\gamma_1=1$ and each later $\gamma_k$ is free to be
chosen arbitrarily large subject to its displayed lower bound. Restricting
these free choices to positive integers is the standard specialization used
for this interpretation; the paper itself does not call them integers.
The book's Series Law identifies $L$ unit resistors in a path with one
resistor of resistance $L$. The distinction between a merely transient tree
and the $\br T>1$ class is located by the original Exercise 3.4, its hint,
the spherical criterion and the growth statement reproduced below.
These background roles are not a newly constructed example of $T$ or a
new proof. No finite-mean or positive-branching-number condition is added
to Exercise 5.57.

The numerical and event expressions in the paper are retained literally:
its induction prints $\gamma_k>2N_{k-1}D_{k-1}\gamma_{k-1}$, while its later
estimate prints $\PP(B_k)\leq N_kD_k\gamma_k/\gamma_{k+1}\leq2^{-k}$.
The displayed unindexed $\cup$ in $B_k$ is also retained. No coefficient,
event or intended erratum has been inserted. Public redistribution permission
for the edited book/paper selection has not been established.
% END CONTEXT 5.57
\normalsize

\subsection*{Original book setup and primary question}
\src{Exercise 5.55, physical page 190, printed page 173; uncounted setup}
% BEGIN SUPPORT SETUP LP:5.55
\noindent\textbf{5.55.}
Randomly stretch a tree $T$ by adding vertices to (``subdividing'') the edges
to make the edge $e(x)$ a path of length $L_x(\omega)$ with $L_x$ i.i.d.
Call the resulting tree $T(\omega)$. Calculate the distribution of
$\br T(\omega)$ in terms of $\br T$ and of the distribution of $L_x$.
Show that if $\br T>1$, then $\br T(\omega)>1$ a.s.
% END SUPPORT SETUP LP:5.55

\src{Exercise 5.57, physical page 190, printed page 173}
% BEGIN PRIMARY STATEMENT 5.57
\noindent\textbf{5.57.}
Consider the stretched tree $T(\omega)$ from Exercise 5.55. If we assume
only that simple random walk on $T$ is transient, must simple random walk
on $T(\omega)$ be transient a.s.?
% END PRIMARY STATEMENT 5.57

\src{Complete Comment 5.57, physical page 631, printed page 614}
% BEGIN PRIMARY COMMENT 5.57
\noindent\textbf{5.57.} See Pemantle and Peres (1996).
% END PRIMARY COMMENT 5.57

\subsection*{Uncounted same-book network and tree background}
\src{Series Law and its printed proof, physical page 44, printed page 27}
% BEGIN SUPPORT SERIES LAW LP
\noindent\textbf{I. Series Law.}
Two resistors\footnote{A resistor $r$ is an edge with resistance $r$.
We have drawn such edges using the squiggly notation common to physics,
but this only indicates that they have weights.}
$r_1$ and $r_2$ in series are equivalent to a single resistor $r_1+r_2$.
In other words, if $w\in\mathsf V(G)\setminus(A\cup Z)$ is a node of degree 2
with neighbors $u_1,u_2$ and we replace the edges $(u_i,w)$ by a single edge
$(u_1,u_2)$ having resistance $r(u_1,w)+r(w,u_2)$, then all potentials and
currents in $G\setminus\{w\}$ are unchanged and the current that flows from
$u_1$ to $u_2$ equals $i(u_1,w)$.
\[
u_1\quad\text{---}\quad w\quad\text{---}\quad u_2
\]
\begin{proof}
It suffices to check that Ohm's and Kirchhoff's laws are satisfied on the
new network for the voltages and currents given. This is easy.
\end{proof}
% END SUPPORT SERIES LAW LP

\src{Exercise 2.8, physical page 47, printed page 30; repeated on physical page 75, printed page 58}
% BEGIN SUPPORT SPHERICAL CRITERION LP:2.8
\noindent\textbf{Exercise 2.8.}
Let $(G,c)$ be a network. A network automorphism of $(G,c)$ is a map
$\phi:G\to G$ that is a bijection of the vertex set with itself and a
bijection of the edge set with itself such that if $x$ and $e$ are incident,
then so are $\phi(x)$ and $\phi(e)$ and such that $c(e)=c(\phi(e))$ for all edges $e$.
Suppose that $(G,c)$ is spherically symmetric about $o$, meaning that if $x$
and $y$ are any two vertices at the same distance from $o$, then there is
an automorphism of $(G,c)$ that leaves $o$ fixed and that takes $x$ to $y$.
Let $C_n$ be the sum of $c(e)$ over all edges $e$ with $d(e^-,o)=n-1$
and $d(e^+,o)=n$. Show that
\[
\mathcal R(o\leftrightarrow\infty)=\sum_{n\geq1}\frac1{C_n},
\]
whence the network random walk on $G$ is transient iff
\[
\sum_{n\geq1}\frac1{C_n}<\infty.
\]
% END SUPPORT SPHERICAL CRITERION LP:2.8

\src{Section 3.3, physical page 99, printed page 82; lower-growth definition excerpt}
In this section, we again consider only locally finite, infinite trees.
To understand better their branching number, we will look at the simpler
notion of growth. For a vertex $x\in T$, let $|x|$ denote its distance from $o$.
Define the lower (exponential) growth rate of a tree $T$ by
\[
\grlow T:=\liminf_{n\to\infty}|T_n|^{1/n},
\]
where $T_n:=\{x\in T;\ |x|=n\}$ is level $n$ of $T$.

\src{Section 3.3, physical page 100, printed page 83; selected paragraph}
We call $T$ spherically symmetric if $\deg x$ depends only on $|x|$.
Recall from Exercise 1.2 that $\br T=\grlow T$ when $T$ is spherically symmetric.

\src{Exercise 3.4, physical page 99, printed page 82; collected again on physical page 109, printed page 92; Comment on physical page 622, printed page 605}
% BEGIN SUPPORT EXERCISE AND HINT LP:3.4
\noindent\textbf{Exercise 3.4.}
For simple random walk on $T$ to be transient, is it necessary that $\br T>1$?

\noindent\textbf{Comment 3.4.} Consider spherically symmetric trees.
% END SUPPORT EXERCISE AND HINT LP:3.4

\clearpage
\section*{Separate author core: Pemantle and Peres}
\src{Undated author copy, physical page 2, printed page 1; complete introductory mathematical text}
% BEGIN SEPARATE AUTHOR CONTEXT PP
Simple random walks on infinite graphs have been widely studied with
particular attention to the dichotomy between recurrence and transience.
In the last two decades there has been interest in random walks in random
environments (RWRE's), which are mixtures of Markov chains in which the
transition probabilities are picked according to some prior distribution.
For reversible Markov chains, the transition probabilities may be described
by associating resistances $R(vw)$ to each edge $(vw)$, the transition
probabilities out of any vertex being proportional to the inverses of the
resistances on the edges incident to that vertex:
\begin{equation}\tag{1}
p(v,w)=\frac{R(vw)^{-1}}{\sum_z R(vz)^{-1}}
\end{equation}
where the sum is over all vertices $z$ adjacent to $v$.
One natural model\footnote[4]{see Remark 1 for another RWRE model relevant to the present discussion},
considered in Grimmett and Kesten (1984), involves taking the resistances
as i.i.d. positive random variables. A result of Adams and Lyons (1991)
implies that if a tree $\Gamma$ has positive Hausdorff dimension, then for
an arbitrary distribution of the resistances the resulting network is
almost surely transient. (See Lyons (1990) where the term branching number
is used for the exponential of the Hausdorff dimension.) The surprising
aspect of this is that no conditions are imposed on the tail of the
distribution. This note proves a converse of this implication, but on a
general graph. Specifically, the implication $(a)\Rightarrow(b)$ in part
$(ii)$ of the theorem below shows that almost sure transience for an
arbitrary resistance distribution on a tree $\Gamma$ implies $p_c<1$ for
IID bond percolation. By results in Lyons (1990), this implies that $\Gamma$
has positive dimension. Part $(i)$ is immediate from the well known
connections between random walks and electrical networks (c.f. Doyle and
Snell (1984)) and is included for completeness.
% END SEPARATE AUTHOR CONTEXT PP

\src{Complete Theorem, physical pages 2--3, printed pages 1--2}
% BEGIN SEPARATE AUTHOR THEOREM PP
\noindent\textbf{Theorem:}
\emph{Let $G$ be a locally finite connected graph. Suppose that a positive
random variable $R(e)$ is attached to every edge $e$ of $G$ so that the
collection $\{R(e)\}$ is i.i.d.. Let $\mu$ denote their common distribution.
Consider the RWRE given by (1).}
\begin{enumerate}
\item[$(i)$] \emph{If $\int x\,d\mu<\infty$ and simple random walk on $G$
is transient, then the resistor network $\{R(e)\}$ has finite net resistance
almost surely, thus defining almost surely a transient random walk.}
\item[$(ii)$] \emph{The following conditions on $G$ are equivalent:
$(a)$ for any distribution $\mu$ on $(0,\infty)$ the resistor network
$\{R(e)\}$ is almost surely transient; $(b)$ there exists some $p<1$ such
that independent bond percolation on $G$ with parameter $p$ contains with
positive probability a cluster on which simple random walk is transient.}
\end{enumerate}
% END SEPARATE AUTHOR THEOREM PP

\src{Complete printed proof, physical pages 3--5, printed pages 2--4}
% BEGIN PRIMARY PROOF CORE 5.57
% BEGIN SEPARATE AUTHOR PROOF PP
\begin{proof}
We rely on two facts from the theory relating reversible random walks to
the corresponding resistor networks. Firstly, the random walk is transient
if and only if the network has finite resistance from some (hence every)
vertex to infinity, i.e., there is some nonzero flow $\{F(e)\}$, satisfying
Kirchoff's laws except at one source vertex, such that the energy
\[
\sum_e F(e)^2R(e)
\]
is finite. Secondly, Rayleigh's principle asserts that the resistance of a
network is monotone in each of the individual resistances. Both facts may
be found in Doyle and Snell (1984). The proof of $(i)$ follows immediately:
$G$ is transient for simple random walk, so some flow $\{F(e)\}$ on $G$ has
finite energy
\[
\sum_e F(e)^2
\]
with respect to unit resistances; the expected energy of this flow in the
network of random resistances is
\[
\EE\sum_e R(e)F(e)^2=\int x\,d\mu\sum_e F(e)^2<\infty
\]
so the resistance is almost surely finite and hence the network is almost surely transient.

$(ii)$ The direction $(b)\Rightarrow(a)$ is easy: assume $(b)$ for some
parameter $p$ and consider the random subgraph of $G$ consisting of those
edges $e$ for which $R(e)\leq Q_p$, where $1>\mu(0,Q_p)\geq p$.
Resistances have the following linearity: multiplying every resistor by $C$
multiplies the resistance of the network by $C$. Thus our assumption $(b)$
implies there is a positive probability that some component of this
subgraph will have finite resistances when all edges have resistance $Q_p$.
By Rayleigh's law, decreasing the resistances on edges in this component
from $Q_p$ to $R(e)$ and decreasing all other resistances from $\infty$ to
$R(e)$ gives a network with finite resistance, and therefore random walk on
$G$ with resistances $\{R(e)\}$ is transient with positive probability.
Since this is a tail event for the countable collection $\{R(e)\}$, the
probability is one.

Now we verify that $(a)\Rightarrow(b)$. Assume that for every $p<1$,
simple random walk on $p$-percolation clusters is almost surely recurrent.
Let $p_k\uparrow1$ and construct a probability measure $\mu$ on
$(0,\infty)$ as follows. Fix a vertex $v\in G$ and set $\gamma_1=1$.
Define a measure $\mu_1$ on $\mathbf R\cup\{\infty\}$ by
$\mu_1(\{\gamma_1\})=p_1$ and $\mu_1(\{\infty\})=1-p_1$.
The RWRE on $G$ with resistance distribution $\mu_1$ is just simple random
walk on the $p_1$-percolation clusters and is thus recurrent. Therefore
there exists $N_1$ such that when choosing a random environment according
to $\mu_1$, the probability is at least $1/2$ that either all edges
adjacent to $v$ have infinite resistance (i.e. $v$ is isolated) or else the
resulting RWRE started from $v$ returns to $v$ in the first $N_1$ steps.
Let $D_1$ be the maximal degree of any vertex in $G$ within distance $N_1$ of $v$.

For the induction step, let $k\geq2$ and assume that $\gamma_j,\mu_j,N_j$
and $D_j$ have been defined for $1\leq j<k$ and $\mu_j(\infty)=1-p_j$.
Pick $\gamma_k>2N_{k-1}D_{k-1}\gamma_{k-1}$. Define $\mu_k$ by moving mass
$p_k-p_{k-1}$ from infinity to $\gamma_k$, i.e.
\[
\mu_k=\mu_{k-1}+(p_k-p_{k-1})(\delta_{\gamma_k}-\delta_\infty).
\]
The RWRE with resistance distribution $\mu_k$ is recurrent since it yields
an environment whose resistances on the $p_k$-percolation cluster are
bounded between 1 and $\gamma_k$ [use the aforementioned linearity together
with Rayleigh's law]. Thus we may choose an integer $N_k$, such that if a
random environment is chosen according to $\mu_k$, the probability is at
most $2^{-k}$ that $v$ is not isolated and the resulting RWRE started from
$v$ fails to return to $v$ within the first $N_k$ steps. Finally, define
$D_k$ to be the maximal degree of any vertex of $G$ within distance $N_k$ of $v$.

Finally, set $\mu(\{\gamma_k\})=p_k-p_{k-1}$ for all $k\geq2$, with
$\mu(\{\gamma_1\})=p_1$, so that $\mu$ is the weak limit of the $\mu_k$.
We construct a $\mu$-RWRE $\{S_n\}$ on the same probability space as a
sequence $\bigl(\{S_n^{(k)}\}:k=1,2,3,\ldots\bigr)$ of $\mu_k$-RWRE's
as follows. Let the resistances $\{R(e)\}$ be i.i.d. with common distribution
$\mu$ and define
\[
R^{(k)}(e)=\begin{cases}R(e)&\text{if }R(e)\leq\gamma_k\\
\infty&\text{if }R(e)>\gamma_k.\end{cases}
\]
Observe that this makes $\{R^{(k)}(e)\}$ i.i.d. with distribution $\mu_k$.
Let $\{S_n\}$ be a random walk starting from $v$ with transition
probabilities determined by the resistances $\{R(e)\}$. Note that if
$R(wz)\leq\gamma_k$ then the transition probability from $w$ to $z$ with
resistances $\{R^{(k)}(e)\}$ is greater than or equal to the the transition
probability from $w$ to $z$ with resistances $\{R(e)\}$. Thus for each $k$,
we may define a random sequence $\{S_n^{(k)}\}$ so as to have transition
probabilities determined by $\{R^{(k)}\}$ and so that $S_n^{(k)}=S_n$ for
all $n\leq T_k$, where $T_k$ is the least time $t$ for which
$R(S_tS_{t+1})>\gamma_k$.

Define events
\begin{align*}
A_k&=\{R(vw)>\gamma_k\text{ for all }w\text{ neighboring }v\}\subseteq\{T_k=0\}\\
B_k&=\{R(S_j,S_{j+1})\geq\gamma_{k+1}\text{ for some }0\leq j<N_k\}\setminus A_k\\
   &=\cup\{0\leq T_k<N_k\}\setminus A_k\\
C_k&=\{S_n^{(k)}\neq v\text{ for all }1\leq n\leq N_k\}.
\end{align*}
The event $G_k=\{S_n\neq v\text{ for all }1\leq n\leq N_k\}$ is contained
in $A_k\cup B_k\cup C_k$, since on $G_k\setminus C_k$ the time $T_k$ is
less than $N_k$. For $1\leq j<N_k$, the probability of the event
$\{T_k=j\}$ is at most $D_k\gamma_k/\gamma_{k+1}$, since the sum of
$R(e)^{-1}$ over edges incident to $S_j$ with resistance greater than
$\gamma_k$ is at most $D_k/\gamma_{k+1}$ and there is at least one edge
incident to $S_j$ with $R(e)^{-1}\geq\gamma_k^{-1}$.
Similarly, $\PP(\{T_k=0\}\setminus A_k)\leq D_k\gamma_k/\gamma_{k+1}$.
Thus $\PP(B_k)\leq N_kD_k\gamma_k/\gamma_{k+1}\leq2^{-k}$ by construction
of $\gamma_{k+1}$. But $N_k$ is defined so that
$\PP(C_k\setminus A_k)\leq2^{-k}$ and clearly $\PP(A_k)\leq1-p_k$.
Therefore
\[
\PP(G_k)\leq(1-p_k)+2^{1-k}\longrightarrow0
\]
as $k\to\infty$ and the $\mu$-RWRE is recurrent.
\end{proof}
% END SEPARATE AUTHOR PROOF PP
% END PRIMARY PROOF CORE 5.57

\src{Remarks 1--2 and acknowledgement, physical page 6, printed page 5}
% BEGIN SEPARATE AUTHOR REMARKS PP
\noindent\textbf{Remarks:}
\begin{enumerate}
\item The recent paper of Grimmett, Kesten and Zhang (1991) shows that for
Euclidean lattices $\mathbf Z^d$ $(d\geq3)$, random walk on supercritical
percolation clusters is almost surely transient, so the theorem is
applicable, and all RWRE's with i.i.d. positive resistances are transient.
Note that the RWRE models discussed in Durrett (1986) and Sunyach (1987)
may be recurrent even when $d\geq3$. Resistances in Durrett's model are
stationary under $\mathbf Z^d$-shifts, so his model is close to the present one.
\item By Lyons (1990), the critical probability for bond percolation on a
tree $\Gamma$ is $p_c(\Gamma)=e^{-\dim\Gamma}$ and if $\dim(\Gamma)>0$
and $p>p_c$, then simple random walk on any infinite cluster is transient
(since the cluster must have positive dimension). For a tree, property
$(b)$ is thus equivalent to $\dim(\Gamma)>0$. The Max Flow Min Cut Theorem
(as applied in Lyons 1990) allows us to formulate this purely in terms of flows:
\begin{quote}
On a tree, property $(a)$ is equivalent to the existence of a flow $\{F(e)\}$
satisfying Kirchoff's laws (except at the root) for unit resistors, and
having the exponential decay property:
\[
F(e)\leq C\rho^{\operatorname{dist}(e)}
\]
where $\operatorname{dist}(e)$ is the distance from $e$ to the root of $\Gamma$.
\end{quote}
\end{enumerate}
\noindent\textbf{Acknowledgement:} Thanks to Harry Kesten for showing us a
surprising example which led to this note.
% END SEPARATE AUTHOR REMARKS PP

\src{Original references, physical pages 6--7, printed pages 5--6}
% BEGIN SEPARATE AUTHOR REFERENCES PP
\begin{thebibliography}{7}
\bibitem{PP1} Adams, S. and Lyons, R. (1991). Amenability, Khazdan's property
and percolation for trees, groups and equivalence relations.
\emph{Israel J. Math.} \textbf{75} 341--370.
\bibitem{PP2} Doyle, P. and Snell, J. L. (1984).
\emph{Random walks and electrical networks}.
Mathematical Association of America: Washington.
\bibitem{PP3} Durrett, R. (1986). Multidimensional random walks in random
environments with subclassical limiting behavior.
\emph{Comm. Math. Phys.} \textbf{104} 87--102.
\bibitem{PP4} Grimmett, G. and Kesten, H. (1984). Random electrical networks
on complete graphs. \emph{J. London Math. Soc.} \textbf{2} 171--192.
\bibitem{PP5} Grimmett, G., Kesten, H. and Zhang, Y. (1991).
Random walk on the infinite cluster of the percolation model. \emph{Preprint}.
\bibitem{PP6} Lyons, R. (1990). Random walks and percolation on trees.
\emph{Ann. Probab.} \textbf{18} 931--958.
\bibitem{PP7} Sunyach, C. (1987). Sur la transience et la r\'ecurrence des
marches al\'eatoires en milieu al\'eatoire.
\emph{Ann. Inst. H. Poin.} \textbf{23} 613--626.
\end{thebibliography}
% END SEPARATE AUTHOR REFERENCES PP
\end{document}
